\textbf{Regularisation.} EWC \cite{kirkpatrick2017overcoming} approximates the posterior over parameters of an old task by a Gaussian with diagonal Fisher precision. Synaptic intelligence \cite{zenke2017continual} is a path-based alternative, learning without forgetting \cite{li2018learning} distils the previous model's outputs on new-task data, and Liu et al.\ use a reparametrisation that approximately diagonalises the Fisher matrix so that EWC's diagonal assumption fits better \cite{liu2018rotate}. We only reimplement EWC.

\textbf{Memory.} Gradient episodic memory \cite{lopezpaz2017gradient} introduced an episodic-memory model together with metrics that include knowledge transfer, and A-GEM \cite{chaudhry2018efficient} is its averaged, cheaper variant. iCaRL \cite{rebuffi2017icarl} uses exemplars with a nearest-mean classifier for class-incremental learning, and DER++ \cite{buzzega2020dark} replays stored logits as well as labels. Chaudhry et al.\ \cite{chaudhry2019tiny} find that a simple baseline that trains jointly on current and stored examples outperforms specifically designed approaches when each example is seen once; our replay baseline follows this plain scheme but, unlike their single-pass setting, trains 30 epochs per task.

\textbf{Evaluation protocol.} van de Ven and Tolias \cite{ven2019three} define three scenarios by whether task identity is given or must be inferred and compare methods on split and permuted MNIST, finding large differences between scenarios. Our split\_til, split\_cil and perm\_dil tasks are the task-, class- and domain-incremental scenarios of that taxonomy on digits instead of MNIST. Our contribution relative to these papers is only the matched small-scale re-measurement and the direct sweep of $M$ against $\lambda$; we make no claim that results transfer to larger data or networks.
